% ============================ CONCLUSION GÉNÉRALE ====================
\chapter*{Conclusion générale}
\addcontentsline{toc}{chapter}{Conclusion générale}
\markboth{Conclusion générale}{Conclusion générale}

Ce stage de Projet de Fin d'Année avait pour ambition de rendre interrogeable en langage
naturel un corpus de documentation technique, \emph{y compris son contenu visuel}, et de
restituer des réponses fiables et vérifiables. Cet objectif est atteint : l'assistant RAG
multimodal développé couvre l'intégralité de la chaîne, du téléversement d'un PDF brut à la
réponse citée illustrée par la figure pertinente.

Sur le plan technique, le projet démontre la maîtrise d'un spectre large et actuel :
analyse documentaire par apprentissage profond (Docling, TableFormer), découpage sémantique
contextualisé, description de figures par modèle de vision-langage (Gemma-3 27B),
représentations vectorielles et bases vectorielles (mistral-embed, ChromaDB), recherche
hybride dense/lexicale fusionnée par RRF, ingénierie de prompts pour la génération
contrainte et le routage (Mistral), développement d'API REST (FastAPI, Pydantic),
authentification moderne (Clerk, JWT/JWKS), interface web réactive (React, TypeScript,
Vite) et évaluation quantitative de systèmes RAG (RAGAS, juge LLM indépendant). Les
résultats de l'évaluation (rappel de contexte parfait, pertinence des réponses de 0{,}93)
valident les choix d'architecture.

Sur le plan de l'ingénierie, la leçon la plus formatrice du stage tient peut-être à son
environnement : développer un système dépendant de modèles d'IA distants sur un poste
verrouillé, derrière un proxy d'inspection TLS, a imposé de diagnostiquer des pannes
réseau opaques, d'arbitrer des changements de fournisseurs en cours de projet (bascule
d'embeddings avec réindexation complète), et de concevoir chaque composant pour la
résilience et la commutabilité. Ces compétences, comprendre \emph{pourquoi} une
infrastructure échoue et concevoir \emph{avec} les contraintes plutôt que contre elles,
sont précisément celles qu'exige l'industrialisation de l'IA en entreprise.

Enfin, ce travail livre à ALTEN un démonstrateur complet et documenté, dont l'architecture
sobre (stockage embarqué, services managés, configuration par environnement) est
directement transposable sur Oracle Cloud Infrastructure.

% ============================ PERSPECTIVES ===========================
\chapter*{Perspectives}
\addcontentsline{toc}{chapter}{Perspectives}
\markboth{Perspectives}{Perspectives}

Plusieurs prolongements, identifiés pendant le stage et volontairement hors périmètre,
constituent des suites naturelles :

\begin{itemize}
  \item \textbf{Déploiement OCI complet} : conteneurisation du backend et du frontend sur
        une VM OCI, cache basculé sur OCI Cache, stockage des fichiers sur Object Storage,
        et éventuellement migration des métadonnées vers Oracle Autonomous Database
        (palier Always Free) ;
  \item \textbf{Re-ranking} : l'ajout d'un ré-ordonnanceur (cross-encoder ou LLM) sur les
        candidats du retrieval hybride est la réponse directe aux scores de
        \emph{context precision} observés sur les corpus à vocabulaire proche ;
  \item \textbf{OCR des documents scannés} : activation de l'option OCR de Docling ou
        recours à un service managé (Azure Document Intelligence), pour couvrir les
        documents-images ;
  \item \textbf{Captioning orienté signification} : itération sur le prompt du VLM pour
        distinguer la description visuelle d'une figure de l'explication de sa
        signification technique, l'axe identifié par le cas d'échec de faithfulness ;
  \item \textbf{Jeu d'évaluation étendu et continu} : constitution d'un jeu de questions
        de référence plus large (génération assistée puis validation humaine) et
        intégration de l'évaluation RAGAS en contrôle continu de non-régression ;
  \item \textbf{Streaming des réponses} : restitution token par token (SSE) pour améliorer
        la latence perçue ;
  \item \textbf{Embeddings multimodaux natifs} : à plus long terme, une voie ColPali
        (late interaction sur pages entières) pourrait compléter ou remplacer le
        captioning pour les corpus où la mise en page elle-même est porteuse de sens ;
  \item \textbf{Gestion des droits documentaires} : cloisonnement du corpus par
        utilisateur ou par équipe, prérequis d'un déploiement au-delà du démonstrateur.
\end{itemize}
